\section{Guestbook Deployment to OpenShift}

\subsection{Configuration Management}

Getting the application running locally is one thing. Deploying it to a real cluster
means thinking carefully about where configuration lives and how secrets are handled.
The rule I followed throughout was simple: nothing sensitive goes into a container image
or into a Git repository. Passwords, tokens, and keys belong in Kubernetes Secrets.
Everything else, hostnames, port numbers, database names, can go into ConfigMaps.

The backend ConfigMap holds the non-sensitive connection details:

\begin{lstlisting}[language=yaml]
apiVersion: v1
kind: ConfigMap
metadata:
  name: backend-config
data:
  DB_HOST: "postgres"
  DB_PORT: "5432"
  DB_USER: "guestbook"
  DB_NAME: "guestbook"
  REDIS_HOST: "redis"
  REDIS_PORT: "6379"
  PORT: "8080"
\end{lstlisting}

The passwords are kept in a Secret. Secrets are not committed to the repository and
must be created directly in the cluster before anything else is applied. The quickest
way to do that is with the imperative command:

\begin{lstlisting}[language=bash]
oc create secret generic backend-secret \
  --from-literal=DB_PASSWORD='password' \
  --from-literal=REDIS_PASSWORD='redispassword'
\end{lstlisting}

The same pattern applies for PostgreSQL and Redis. Their secrets are created the same
way before the StatefulSets are applied:

\begin{lstlisting}[language=bash]
oc create secret generic postgres-secret \
  --from-literal=POSTGRESQL_PASSWORD='password'

oc create secret generic redis-secret \
  --from-literal=REDIS_PASSWORD='redispassword'
\end{lstlisting}

\subsection{Resource Manifests}

\subsubsection{Backend}

The backend is deployed as a standard Deployment with two replicas. It pulls the
ConfigMap values as a block using \texttt{envFrom} and then adds the password values
individually from the Secret using \texttt{secretKeyRef}:

\begin{lstlisting}[language=yaml]
apiVersion: apps/v1
kind: Deployment
metadata:
  name: backend
  labels:
    app: backend
spec:
  replicas: 2
  selector:
    matchLabels:
      app: backend
  template:
    metadata:
      labels:
        app: backend
    spec:
      containers:
        - name: backend
          image: ghcr.io/sirajwahaj/ocp-guestbook-backend:sha-1b8b5d1
          imagePullPolicy: Always
          ports:
            - name: http
              containerPort: 8080
          envFrom:
            - configMapRef:
                name: backend-config
          env:
            - name: DB_PASSWORD
              valueFrom:
                secretKeyRef:
                  name: backend-secret
                  key: DB_PASSWORD
            - name: REDIS_PASSWORD
              valueFrom:
                secretKeyRef:
                  name: backend-secret
                  key: REDIS_PASSWORD
\end{lstlisting}

The Service for the backend is straightforward and just routes traffic on port 8080 to
whichever pods carry the \texttt{app: backend} label.

\subsubsection{Frontend}

The frontend Deployment is similar but simpler since it does not need any environment
variables. It is also exposed externally through an OpenShift Route:

\begin{lstlisting}[language=yaml]
apiVersion: route.openshift.io/v1
kind: Route
metadata:
  name: frontend
spec:
  to:
    kind: Service
    name: frontend
  port:
    targetPort: 8080
  tls:
    termination: edge
\end{lstlisting}

I intentionally left the \texttt{host} field out of the Route. When you set a hardcoded
hostname that does not match the cluster domain, OpenShift just serves a placeholder
page and you spend a while wondering why it is not working. Letting OpenShift generate
the hostname automatically avoids that entirely.

\subsubsection{PostgreSQL}

PostgreSQL is deployed as a StatefulSet because it needs a stable identity and persistent
storage. A regular Deployment would not guarantee that the same pod gets the same volume
after a reschedule:

\begin{lstlisting}[language=yaml]
apiVersion: apps/v1
kind: StatefulSet
metadata:
  name: postgres
spec:
  selector:
    matchLabels:
      app: postgres
  replicas: 1
  template:
    metadata:
      labels:
        app: postgres
    spec:
      containers:
        - name: postgres
          image: quay.io/fedora/postgresql-16:latest
          ports:
            - containerPort: 5432
          volumeMounts:
            - name: postgres
              mountPath: /var/lib/pgsql/data
          env:
            - name: POSTGRESQL_USER
              valueFrom:
                configMapKeyRef:
                  name: postgres-config
                  key: POSTGRESQL_USER
            - name: POSTGRESQL_PASSWORD
              valueFrom:
                secretKeyRef:
                  name: postgres-secret
                  key: POSTGRESQL_PASSWORD
            - name: POSTGRESQL_DATABASE
              valueFrom:
                configMapKeyRef:
                  name: postgres-config
                  key: POSTGRESQL_DATABASE
  volumeClaimTemplates:
    - metadata:
        name: postgres
      spec:
        accessModes: ["ReadWriteOnce"]
        storageClassName: "gp3"
        resources:
          requests:
            storage: 1Gi
\end{lstlisting}

\subsubsection{Redis}

Redis is also a StatefulSet, but without persistent storage since it only acts as a
cache. If it restarts and loses its data, the backend will simply fall back to PostgreSQL
on the next request and repopulate the cache. The manifest is the same structure as
PostgreSQL but without \texttt{volumeClaimTemplates}.

\subsection{Problems Encountered During Deployment}

\subsubsection{Backend Pods in CrashLoopBackOff}

The backend pods kept restarting when I first applied the manifests. Running
\texttt{oc logs} showed that the backend was failing to connect to PostgreSQL and Redis,
which had not finished starting yet. The fix was to make sure the ConfigMaps and Secrets
existed before applying the Deployment, and to delete the old pods so they could restart
cleanly with the correct environment variables.

\subsubsection{PostgreSQL Pod Stuck in Pending}

The PostgreSQL pod never moved out of Pending state. The event log showed that the
StorageClass referenced in the manifest did not exist on this cluster. I had copied a
manifest that used \texttt{nfs-storage} but the Sandbox uses \texttt{gp3}. Updating
\texttt{storageClassName} in the \texttt{volumeClaimTemplates} and reapplying fixed it.

\subsubsection{Immutable Fields on StatefulSet}

After changing the StorageClass I tried to apply the updated StatefulSet directly, which
failed with a forbidden error. Fields like \texttt{volumeClaimTemplates} and
\texttt{storageClassName} cannot be changed on an existing StatefulSet. The only options
are to delete the StatefulSet first and then reapply, or to only update fields that are
allowed to change, like \texttt{replicas} or the container image. I deleted it and
reapplied, accepting that the persistent volume would also be lost in this case since I
was still in testing.

\subsubsection{Frontend Not Reachable Externally}

Even after the Route was created the frontend was not reachable from a browser. The
Route had a hardcoded \texttt{host} field that did not match the cluster domain, so
requests were reaching the wrong place. Removing that field let OpenShift assign the
correct hostname automatically. I also added a \texttt{readinessProbe} to the frontend
Deployment so the Router would not forward traffic until the pod was actually ready to
serve requests.
